\subsection{Essentially finitely generated algebras}

Let \(k\) a ring. Let A be an \(k\) -algebra and \(\mathbf{x} = (x_i)_{i \in I}\) a family of elements of A; denote by \(\mathbf{A}'\) the subalgebra of A generated by the \(x_i\) . We say that \(\mathbf{x}\) is an essentially generating family of the \(k\) -algebra A if, for every element \(a\) of A, there exists an element \(s\) of \(\mathbf{A}'\) , invertible in A, such that \(sa\) belong to \(\mathbf{A}'\) . This is equivalent to saying that, for every \(a \in \mathbf{A}\) , there exist polynomials P and Q in \(k[(\mathrm{X}_i)_{i \in I}]\) such that \(\mathrm{Q}(\mathbf{x})\) is invertible in A and that one has \(a = \mathrm{P}(\mathbf{x})\mathrm{Q}(\mathbf{x})^{-1}\) .

We shall say that a \(k\) -algebra A is essentially of finite type if it admits a finite essentially generating family. This is equivalent to saying that there exists a \(k\) -algebra A' of finite type and a multiplicative subset S of A' such that the \(k\) -algebra A is isomorphic to \(\mathrm{S}^{-1}\mathrm{A}'\) .

Examples.--- 1) Saying that an extension L of a field K is an essentially finitely generated K-algebra means that it is a finitely generated extension in the sense of A, V, p.~11, def. 2. The K-algebra L is of finite type only if it is a finite-degree extension of K (V, § 3, no. 4, cor. 3 of th. 3).

\begin{enumerate}
\def\labelenumi{\arabic{enumi})}
\setcounter{enumi}{1}
\tightlist
\item
  For a \(k\) -algebra to be local and essentially of finite type, it is necessary and sufficient that it be isomorphic to a \(k\) -algebra of the form \(\mathsf{A}_{\mathfrak{p}}\) , where A is a \(k\) finitely generated algebra and \(\mathfrak{p}\) a prime ideal of A (cf.~II, § 2, no. 5, prop. 11 (iii)).
\end{enumerate}

PROPOSITION 1. - If the ring k is Noetherian, every essentially finitely generated k-algebra is a Noetherian ring.

This follows from III, § 2, no. 10, cor. 3 of th. 2 and II, § 2, no. 4, cor. 2 of prop. 10.

From the definition, we deduce the following properties:

PROPOSITION 2.--- a) Every quotient algebra of a k-algebra essentially of finite type is essentially of finite type.

\begin{enumerate}
\def\labelenumi{\alph{enumi})}
\setcounter{enumi}{1}
\item
  Every ring of fractions of an essentially finitely generated k-algebra is an essentially finitely generated k-algebra.
\item
  The \(k\) -algebra product of a finite family of \(k\) -algebras essentially of finite type is essentially of finite type.
\item
  Let \(k \to k'\) a ring homomorphism; for every \(k\) -algebra A essentially of finite type, the \(k'\) -algebra \(A_{(k')} = k' \otimes_k \Lambda\) is essentially of finite type.
\end{enumerate}

COROLLARY.- Let A be an \(k\) -algebra essentially of finite type and B a \(k\) -algebra that is Noetherian. Then the ring \(\mathrm{A} \otimes_k \mathrm{B}\) is Noetherian.

Indeed, it is a B-algebra essentially of finite type (prop. 2, d)) and one applies proposition 1.

PROPOSITION 3. Let \(k\) a ring, A a linearly topologized formally smooth \(k\) -algebra and B an A-algebra. If A is essentially of finite type over \(k\) and B essentially of finite type over A, then B is essentially of finite type over \(k\) .

Let \(\rho: A \to B\) the canonical map. Let \(x = (x_i)_{i \in I}\) a finite essentially generating family of the \(k\) -algebra \(A\) and \(A'\) the subalgebra it generates; let \(y = (y_j)_{j \in J}\) a finite essentially generating family of the \(A\) -algebra \(B\) . Let \(B'\) the sub- \(k\) -algebra of \(B\) generated by the \(\rho(x_i)\) and the \(y_j\) . Let \(b \in B\) ; by hypothesis there exist polynomials \(P\) and \(Q\) in \(A[(Y_j)_{j \in J}]\) such that \(Q(y)\) let it be invertible in \(B\) and that we have \(Q(y)b = P(y)\) The nonzero coefficients of \(P\) and \(Q\) are finite in number; there exists a polynomial \(R \in k[(X_i)_{i \in I}]\) such that \(R(x)\) let it be invertible in \(A\) and that \(R(x)P\) and \(R(x)Q\) belong to \(A'[[(Y_j)_{j \in J}]\) Then \(\rho(R(x))Q(y)\) be invertible in \(B\) and one has

\[
\rho (\mathrm{R} (\mathbf {x})) \mathrm{Q} (\mathbf {y}) b = \rho (\mathrm{R} (\mathbf {x})) \mathrm{P} (\mathbf {y}) \in \mathrm{B} ^ {\prime}.
\]

Thus the elements \(\rho(x_i)\) for \(i \in I\) and \(y_j\) for \(j \in J\) form an essentially finitely generating family of the \(k\) -algebra B.

COROLLARY. -- The tensor product of two k-algebras essentially of finite type is a k-algebra essentially of finite type.

Indeed, let A and B be two essentially finitely generated k-algebras. Then \(A \otimes_{k} B\) is essentially of finite type over A (prop. 2, d)) hence over k (prop. 3).

